\documentclass[10pt]{article}

\usepackage[utf8]{inputenc}

\usepackage[T1]{fontenc}

\usepackage{graphicx}

\usepackage[export]{adjustbox}

\graphicspath{ {./images/} }

\usepackage{amsmath}

\usepackage{amsfonts}

\usepackage{amssymb}

\usepackage[version=4]{mhchem}

\usepackage{stmaryrd}

\usepackage{mathrsfs}

\usepackage{bbold}



\DeclareUnicodeCharacter{0131}{$\imath$}



\begin{document}

\includegraphics[max width=\textwidth, center]{2023_07_08_538c59e7d0ee53142860g-1(4)}

цепь с переходными вероятностями



\[

p_{i j}=\mathrm{P}\left\{X\left(\tau_{n}\right)=j \mid X\left(\tau_{n-1}\right)=i\right\} .

\]



Распределения моментов скачков $\tau_{n}$ описываются с помощью функций распределения $F_{i j}(x)$ следующим образом:



$\mathrm{P}\left\{\tau_{n}-\tau_{n-1} \leqslant x, \quad X\left(\tau_{n}\right)=j \mid X\left(\tau_{n-1}\right)=i\right\}=p_{i j} F_{i j}(x)$ (и при әтом не зависят от состояний процесса в более ранние моменты времени). Если



\[

F_{i j}^{\prime}(x)=e^{-\lambda_{i} j^{x}}, x \geqslant 0,

\]



для всех $i, j \in N$, то П. п. $X(t)$ является цепью Маркова с непрерывным временем. Если все распределения вырождены в одной точке, то получают цепь Маркова с дискретным временем.



П. п. служит моделью многих процессов массового обслуживания и теории надежности. С П. П. связаны шродессы марковского восстановления (см. Восстановления теория), описывающие количество посещений процессом $X(t)$ состояний $i \in N$ за время $[0, t]$.



Изучение П. п. и марковских процессов восстановления аналитически сводится к системе интегральных уравнений восстановления.



Лит.: [1] К о р о люк В. С., Т у р би н А. Ф., Полумарковские процессы и их приложения, К., 1976 .



ПОЛУ МАРТИНГАЛ - понятие, равносильное понятиям субмартингала или супермартингала. Именно, стохастич. последовательность $X=\left(X_{t}, \mathscr{\mathscr { T }} t\right), t \in T \subseteq[0, \infty)$, заданная на вероятностном пространстве $(\Omega, \mathscr{F}, \mathrm{P})$ с выделенным на нем убывающим семейством $\sigma$-алгебр $\left(\mathscr{F}_{t}\right)_{t \in T}, \mathscr{F}_{s} \subseteq \mathcal{F} t \subseteq \mathscr{F}, s \leqslant t$, наз. п о л у м а р т и н $г$ ал о м, если $\mathrm{E}\left|X_{t}\right|<\infty, X_{t}$ является $\mathscr{F}_{t}$-измеримой и с вероятностью 1 или



\[

\begin{aligned}

& \mathrm{E}\left(X_{t} \mid \mathscr{F}_{s}\right) \geqslant X_{s}, \\

& \mathrm{E}\left(X_{t} \mid \mathscr{F}_{s}\right) \leqslant X_{s} .

\end{aligned}

\]



или



В случае (1) П. наз. с у б м а р т и нг а ло м, в случаө (2) - с у п е р м а р т и нг а лом.



В современной литературе термин «П.»или не употребляют, или отождествляют с понятием субмартингала (супермартингал определяется изменением знака из субмартингала и наз. иногда нижним П.). См. также Мартингал. A. В. Прохоров.



ПОЛУНАСЛЕДСТВЕННОЕ КОЛЬЦО с л Р в акольцо, все конечно порожденные левые идеалы к-рого ıроективны. П. к. являются кольцо целых чисел, кольцо многочленов от одного неизвестного над полем, регулярные кольца в смысле Неймана, наследственные кольца, кольца конечно порожденных свободных идеалов (полу-FI-кольца). Аналогично определяется правое П. к. Левое П. к. не обязано быть правым П. к. Однако локальное левое П. к. оказывается областью целостностк и правым П. к. Кольцо матриц над П. к. является П. к. Если $R$ есть П. к. и $e^{2}=e \in R$, то $e R e$ есть П. к. Конечно порожденный подмодуль проективного модуля над П. к. изоморфен прямой сумме нек-рого множества конечно порожденных левых идеалов основного кольца и, следовательно, проективен. Каждый такой модуль может быть представлен и как прямая сумма модулей, двойственных конечно порожденным правым идеалам основного кольца.



Для коммутативного кольца $R$ эквивалентны следующие свойства: (1) $R$ есть П. к.; (2) $(A \cap B) C=A C \cap B C$, где $A, B$ и $C$ - произвольные идеалы кольца $R ;(3)$ полное кольцо частных кольца $R$ регулярно в смысле Неймана, и для всякого максимального идеала $\mathfrak{m}$ кольца $R$ кольцо частных $R_{\mathrm{m}}$ является кольцом нормирования; (4) все 2-порожденные идеалы кольца $R$ проектив- ны. Кольцо многочлөнов от одного переменного над коммутативным кольцом $R$ оказывается П. к. в том и только в том случае, когда $R$ регулярно в смысле Неймана.



Лum.: [1] К а р т а н А., Э йл е н б е р г С., Гомологическая алгебра, пер. с англ., М., 1960; [2] Итоги науки и техни-

ки. Алгебра. Топология. Геометрия, т. 14, М., 1976, с. 57-190, ки. Алгебра. Топология. Геометрия, т. 14, М., 1976, с. 57-190,

т. 19, М., 1981, с. 31-134. А. Скорняков.



ПОЛУНЕПРЕРЫВНАЯ ФУНКЦИЯ - функция из первого Бэра класса. Подробнее, числовая функция $f$, определенная на полном метрич. пространстве $X$, наз. полунепре ры вной снизу (св е р х у в точке $x_{0} \in X$, если



\[

\underline{\lim }_{x \rightarrow x_{0}} f(x)=f\left(x_{0}\right)\left(\overline{\operatorname{im}}_{x \rightarrow x_{0}} f(x)=f\left(x_{0}\right)\right) .

\]



Функция $f$ наз. полунепрерывной снизу (сверху) на $X$, если она полунепрерывна снизу (сверху) для всех $x \in X$. Предел монотонно возрастающей (убывающей) последовательности полунепрерывных снизу (сверху) в точке $x_{0}$ функций есть П. ф. снизу (сверху) в $x_{0}$. Если $u(x)$ и $v(x)$ есть $\Pi$. Ф. соответственно снизу и сверху на $X$ и для всех $x \in X$ имеет место $u(x) \leqslant v(x), u(x)>-\infty, v(x)<$ $<+\infty$, то существует непрерывная на $X$ функция $f$ такая, что $v(x) \leqslant f(x) \leqslant u(x)$ для всех $x \in X$. Если $\mu-$ неотрицательная мера на $\mathbb{R}^{n}$, то для любой $\mu$-измеримой функции $f: \mathbb{R}^{n} \rightarrow \mathbb{R}$ существуют две последовательности функций $\left\{u_{n}(x)\right\}$ и $\left\{v_{n}(x)\right\}$, удовлетворяющие условиям: 1) $u_{n}(x)$ полунепрерывны снизу, $v_{n}(x)$ полунепрерывны сверху, 2) каждая функция $u_{n}(x)$ ограничена снизу, каждая функция $v_{n}(x)-$ сверху, 3) последовательность $\left\{u_{n}\right\}$ невозрастающая, последовательность $\left\{v_{n}\right\}$ неубывающая, 4) для всех $x$ имеет место неравенство



\[

u_{n}(x) \geqslant f(x) \geqslant v_{n}(x),

\]



\begin{center}

\includegraphics[max width=\textwidth]{2023_07_08_538c59e7d0ee53142860g-1(2)}

\end{center}



\[

\lim _{n \rightarrow \infty} u_{n}(x)=\lim _{n \rightarrow \infty} v_{n}(x)=f(x),

\]



\begin{enumerate}

  \setcounter{enumi}{5}

  \item если для $E \subset \mathbb{R} n$ функция $f$ суммируема, $f \in L(E)$, то $u_{n}, v_{n} \in L(E) \quad \mathbf{u}$

\end{enumerate}



\[

\lim _{n \rightarrow \infty} \int_{E} u_{n} d \mu=\lim _{n \rightarrow \infty} \int_{E} v_{n} d \mu=\int_{E} f d \mu

\]



(те о р е ма В и т а ли- К а р а те о дор и).



\includegraphics[max width=\textwidth, center]{2023_07_08_538c59e7d0ee53142860g-1(1)}

венной переменной, 3 изд., М., 1974; [2] С а к с С., Теория

интеграла, пер. с англ., М., 1949.

И. А. Виноградова.



ПОЛУНЕПРЕРЫВНОЕ ОТОБРАЖЕНИЕ С в Р рх у (с н и з у) - отображение $f$ одного топологич. пространства $X$ в другое $Y$ такое, что из



\[

\lim x_{n}=x

\]



следует



\begin{center}

\includegraphics[max width=\textwidth]{2023_07_08_538c59e7d0ee53142860g-1}

\end{center}



(здесь $\overline{\lim }(\underline{\lim )}$ - верхний (нижний) предел).



II0ЛУHEIPEPH, М. И. Войцеховский.



\includegraphics[max width=\textwidth, center]{2023_07_08_538c59e7d0ee53142860g-1(3)}

(с в е р х у) - разбиение $D$, т. е. замкнутое дизъюнктное покрытие пространства $X$, такое, что проекция $\rho: X \rightarrow D$ является открытым (замкнутым) отображением. $M$. И. Войчеховский.



ПОЛУ НЕПРЕРЫВНЫЙ МЕТОД СУММИРОВАНИЯ - метод суммирования рядов и последовательностей, определенный с помощью последовательности функций. Пусть $\left\{a_{k}(\omega)\right\}, k=0,1, \ldots,-$ последовательность функций, заданных на нек-ром множестве $E$ изменения параметра $\omega$, и $\omega_{0}$ - точка сгущения этого множества (конечная или бесконечная). Данную последовательность $\left\{s_{n}\right\}$ с помощью функций $a_{k}(\omega)$ преобразуют в функцию $\sigma(\omega)$ :



\[

\sigma(\omega)=\sum_{k=0}^{\infty} a_{k}(\omega) s_{k}

\]





\end{document}